%%%%%%%%%%%%%%%%%
% This is an sample CV template created using altacv.cls
% (v1.7.2, 28 August 2024) written by LianTze Lim (liantze@gmail.com). Compiles with pdfLaTeX, XeLaTeX and LuaLaTeX.
%
%% It may be distributed and/or modified under the
%% conditions of the LaTeX Project Public License, either version 1.3
%% of this license or (at your option) any later version.
%% The latest version of this license is in
%%    http://www.latex-project.org/lppl.txt
%% and version 1.3 or later is part of all distributions of LaTeX
%% version 2003/12/01 or later.
%%%%%%%%%%%%%%%%

%% Use the "normalphoto" option if you want a normal photo instead of cropped to a circle
% \documentclass[10pt,a4paper,normalphoto]{altacv}

\documentclass[10pt,a4paper,ragged2e,withhyper]{altacv}
\graphicspath{{images/}}
%% AltaCV uses the fontawesome5 and simpleicons packages.
%% See http://texdoc.net/pkg/fontawesome5 and http://texdoc.net/pkg/simpleicons for full list of symbols.

% Change the page layout if you need to
\geometry{left=1.25cm,right=1.25cm,top=1.5cm,bottom=1.5cm,columnsep=1.2cm}

% The paracol package lets you typeset columns of text in parallel
\usepackage{paracol}

% Change the font if you want to, depending on whether
% you're using pdflatex or xelatex/lualatex
% WHEN COMPILING WITH XELATEX PLEASE USE
% xelatex -shell-escape -output-driver="xdvipdfmx -z 0" sample.tex
\ifxetexorluatex
  % If using xelatex or lualatex:
  \setmainfont{Roboto Slab}
  \setsansfont{Lato}
  \renewcommand{\familydefault}{\sfdefault}
\else
  % If using pdflatex:
  \usepackage[rm]{roboto}
  \usepackage[defaultsans]{lato}
  % \usepackage{sourcesanspro}
  \renewcommand{\familydefault}{\sfdefault}
\fi

% Change the colours if you want to
\definecolor{SlateGrey}{HTML}{2E2E2E}
\definecolor{LightGrey}{HTML}{666666}
\definecolor{DarkPastelRed}{HTML}{450808}
\definecolor{PastelRed}{HTML}{8F0D0D}
\definecolor{GoldenEarth}{HTML}{E7D192}
\colorlet{name}{black}
\colorlet{tagline}{PastelRed}
\colorlet{heading}{DarkPastelRed}
\colorlet{headingrule}{GoldenEarth}
\colorlet{subheading}{PastelRed}
\colorlet{accent}{PastelRed}
\colorlet{emphasis}{SlateGrey}
\colorlet{body}{LightGrey}

% Change some fonts, if necessary
\renewcommand{\namefont}{\Huge\rmfamily\bfseries}
\renewcommand{\personalinfofont}{\footnotesize}
\renewcommand{\cvsectionfont}{\LARGE\rmfamily\bfseries}
\renewcommand{\cvsubsectionfont}{\large\bfseries}


% Change the bullets for itemize and rating marker
% for \cvskill if you want to
\renewcommand{\cvItemMarker}{{\small\textbullet}}
\renewcommand{\cvRatingMarker}{\faCircle}
% ...and the markers for the date/location for \cvevent
% \renewcommand{\cvDateMarker}{\faCalendar*[regular]}
% \renewcommand{\cvLocationMarker}{\faMapMarker*}


% If your CV/résumé is in a language other than English,
% then you probably want to change these so that when you
% copy-paste from the PDF or run pdftotext, the location
% and date marker icons for \cvevent will paste as correct
% translations. For example Spanish:
% \renewcommand{\locationname}{Ubicación}
% \renewcommand{\datename}{Fecha}


%% Use (and optionally edit if necessary) this .tex if you
%% want to use an author-year reference style like APA(6)
%% for your publication list
% \input{pubs-authoryear.tex}

%% Use (and optionally edit if necessary) this .tex if you
%% want an originally numerical reference style like IEEE
%% for your publication list
\input{pubs-num.tex}
%% marcolanconelli-cv.bib contains your publications
\addbibresource{marcolanconelli-cv.bib}

\begin{document}
\name{Marco Lanconelli}
\today

\tagline{Software Engineer}
%% You can add multiple photos on the left or right
\photoR{2.8cm}{me}
% \photoL{2.5cm}{Yacht_High,Suitcase_High}

\personalinfo{%
  % Not all of these are required!
  \printinfo{\faBirthdayCake}{09 Lug 1993}
  % \mailaddress{}
  \location{Bologna, ITALY}
  \phone{+39 3342267655}
  \email{marco.lanconelli@outlook.it}
  %\homepage{www.homepage.com}
  % \twitter{@twitterhandle}
  %\xtwitter{@x-handle}
  \linkedin{marco-lanconelli}
  \github{mistahuman}
  %\orcid{0000-0000-0000-0000}
  %% You can add your own arbitrary detail with
  %% \printinfo{symbol}{detail}[optional hyperlink prefix]
  % \printinfo{\faPaw}{Hey ho!}[https://example.com/]

  %% Or you can declare your own field with
  %% \NewInfoFiled{fieldname}{symbol}[optional hyperlink prefix] and use it:
  % \NewInfoField{gitlab}{\faGitlab}[https://gitlab.com/]
  % \gitlab{your_id}
  %%
  %% For services and platforms like Mastodon where there isn't a
  %% straightforward relation between the user ID/nickname and the hyperlink,
  %% you can use \printinfo directly e.g.
  % \printinfo{\faMastodon}{@username@instace}[https://instance.url/@username]
  %% But if you absolutely want to create new dedicated info fields for
  %% such platforms, then use \NewInfoField* with a star:
  % \NewInfoField*{mastodon}{\faMastodon}
  %% then you can use \mastodon, with TWO arguments where the 2nd argument is
  %% the full hyperlink.
  % \mastodon{@username@instance}{https://instance.url/@username}
}

\makecvheader
%% Depending on your tastes, you may want to make fonts of itemize environments slightly smaller
% \AtBeginEnvironment{itemize}{\small}

%% Set the left/right column width ratio to 6:4.
\columnratio{0.6}

% Start a 2-column paracol. Both the left and right columns will automatically
% break across pages if things get too long.
\begin{paracol}{2}
\cvsection{Esperienza}

\cvevent{Software Engineer - Fullstack Developer}{Freelancer}{Ottobre 2025 -- Oggi}{Bologna}
Sviluppo fullstack di una piattaforma di orchestrazione cloud multi-tenant, nell'ambito del programma europeo IPCEI-CIS:
\begin{itemize}
    \item servizi backend in Go conformi agli standard TM Forum Open API;
    \item modello di ownership multi-tenant propagato tra i servizi;
    \item operatori Kubernetes su controller-runtime, con CRD custom per ordini e risorse;
    \item integrazione con provider cloud eterogenei;
    \item console utente in Vuejs 3 e relativa suite di test end-to-end.
\end{itemize}

\medskip
\cvevent{Software Engineer}{NIER Ingegneria S.p.A. - Società Benefit}{Ottobre 2020 -- 2025}{Bologna}
Progettazione e implementazione dell'intero ciclo di sviluppo software safety-critical, inclusi:
\begin{itemize}
    \item design e sviluppo software per sistemi ERTMS (interlocking, RBC);
    \item progettazione di protocolli di comunicazione safety-critical;
    \item sviluppo di tool di simulazione e di generazione automatica delle configurazioni;
    \item automazione dei processi di sviluppo;
    \item verifica e validazione del codice.
\end{itemize}

Sviluppo fullstack di applicazioni web per scopi di innovazione aziendale, incluse piattaforme per la gestione dei processi aziendali e per gestione documentale.


%\divider

%\cvevent{Tirocinio Curriculare}{ALMAPLASMA}{Apr 2015 -- Oct 2015}{Bologna}
%Trattamenti tramite plasmi di non equilibrio in ambito biomedicale.


\medskip
\cvsection{Istruzione}

\cvevent{Laurea Magistrale in Ingegneria Energetica e Nucleare}{Università degli Studi di Bologna}{Set 2016 -- Dic 2019}{Bologna \hfill \faMedal 110L/110L}

\textbf{Modellazione di un reattore nucleare LFR e implementazione su piattaforma numerica con accoppiamento neutronico e termoidraulico} (Relatore: Sandro Manservisi)

%\textit{Il reattore nucleare che è stato integrato all’interno della piattaforma numerica si basa sul progetto ALFRED, il quale si pone lo scopo di dimostrare la fattibilità dei reattori veloci raffreddati al piombo. La piattaforma numerica permette di ottenere una simulazione comprensiva della fisica del reattore, attraverso l’implementazione del codice di neutronica Dragon/Donjon e di termoidraulica FEMuS.}

\divider

% \cvevent{Laurea Triennale in Ingegneria Energetica}{Alma Mater Studiorum, Università di
% Bologna}{Set 2012 -- Mar 2016}{Bologna \hfill \faMedal 95/110L}

% \divider

% \cvevent{Maturità Scientifica}{Liceo Scientifico Tecnologico, G. R. Curbastro}{Set 2007 -- Giu 2012}{Lugo (RA) \hfill \faMedal 80/100}

\cvsection{Pubblicazioni}
\cvevent{}{\printinfo{\faUsers}{Conference Proceedings}}{Set 2019}{Milan}
Marco Lanconelli, Sandro Manservisi, Marco Sumini 

\textbf{"A Multiphysics Approach to LFR Analysis"} \textit{Technical Meeting on the Benefits and Challenges of Fast Reactors of the SMR Type}

%\cvsection{A Day of My Life}

% Adapted from @Jake's answer from http://tex.stackexchange.com/a/82729/226
% \wheelchart{outer radius}{inner radius}{
% comma-separated list of value/text width/color/detail}
%\wheelchart{1.5cm}{0.5cm}{%
%  7/8em/accent!30/{Sleep,\\beautiful sleep},
%  3/8em/accent!40/Playing League of Legends,
%  8/8em/accent!60/Daytime job,
%  1/6em/accent!20/Fapping
%}

% use ONLY \newpage if you want to force a page break for
% ONLY the current column
\newpage



%% Switch to the right column. This will now automatically move to the second
%% page if the content is too long.
\switchcolumn

%\cvsection{My Life Philosophy}
%\begin{quote}
%``Ma scusa io che cazzo ne so.''
%\end{quote}

\cvsection{Progetti}

\cvevent{MatriSkillS}{NIER Ingegneria S.p.A. - Società Benefit}{2021 - 2025}{}
\begin{itemize}
    \item Ruolo di sviluppatore fullstack
\end{itemize}
Piattaforma software per la gestione e la crescita del personale, oltre che per la gestione dei processi aziendali, inclusi selezione dei candidati e ricerca di competenze. \\
\href{https://www.niering.it/software-gestionali-hr-le-nuove-funzionalita-di-matriskills/}{(\underline{dettagli qui)}}

\textbf{\small Docker - Flask - MongoDB - Vue.js - Bootstrap}

\divider

\cvevent{EmuXP}{NIER Ingegneria S.p.A. - Società Benefit}{2022 - 2025}{}
\begin{itemize}
    \item Ruolo di sviluppatore fullstack
\end{itemize}
Suite di simulazione multiprotocollo per la validazione funzionale di nodi di comunicazione, con strumenti avanzati per la manipolazione dei protocolli. \href{https://www.niering.it/case_histories/simulatore-multiprotocollo-emuxp/}{(\underline{dettagli qui)}}

\textbf{\small C - Cython - gRPC - Flask - Vue.js - Vuetify}


%\cvsection{Molto fiero di}
%\cvachievement{\faTrophy}{Platino}{League of Legends}

%\divider

%\cvachievement{\faHeartbeat}{4.8k gearscore}{World of Warcraft}

%\divider

%\cvachievement{\faHeartbeat}{0 gioie}{Solo fallimenti}

\cvsection{Competenze}
\cvtag{Problem solving}\\
\cvtag{Attenzione per i dettagli}

\divider\smallskip

\cvtag{Embedded Systems}
\cvtag{FreeBSD}\\
\cvtag{Linux}
\cvtag{Bash}
\cvtag{C}\\
\cvtag{Javascript}
\cvtag{Python}
\cvtag{Go}\\
\cvtag{Vuejs}
\cvtag{Svelte}
\cvtag{React}
\cvtag{Git}\\
\cvtag{SvelteKit}
\cvtag{Docker}
\cvtag{LaTeX}\\
\cvtag{Kubernetes}
\cvtag{Helm}


\divider\smallskip

\cvtag{PVS}
\cvtag{EURORADIO}
\cvtag{MISRA C}\\
\cvtag{UNISIG}

\cvsection{Lingue}
\cvskill{Italiano}{4.5}
\divider
\cvskill{Inglese}{3}
%\divider
%\cvskill{Francese}{2}
%\divider



%% Yeah I didn't spend too much time making all the
%% spacing consistent... sorry. Use \smallskip, \medskip,
%% \bigskip, \vspace etc to make adjustments.

\medskip

%\cvsection{Pubblicazioni}
%% Specify your last name(s) and first name(s) as given in the .bib to automatically bold your own name in the publications list.
%% One caveat: You need to write \bibnamedelima where there's a space in your name for this to work properly; or write \bibnamedelimi if you use initials in the .bib
%% You can specify multiple names, especially if you have changed your name or if you need to highlight multiple authors.
%\mynames{Lim/Lian\bibnamedelima Tze,
%  Wong/Lian\bibnamedelima Tze,
%  Lim/Tracy,
%  Lim/L.\bibnamedelimi T.}
%% MAKE SURE THERE IS NO SPACE AFTER THE FINAL NAME IN YOUR \mynames LIST

%\nocite{*}

%\printbibliography[heading=pubtype,title={\printinfo{\faBook}{Books}},type=book]

%\divider

%\printbibliography[heading=pubtype,title={\printinfo{\faFile*[regular]}{Journal Articles}},type=article]

%\divider

%\printbibliography[heading=pubtype,title={\printinfo{\faUsers}{Conference Proceedings}},type=inproceedings]



%\cvsection{Referees}

% \cvref{name}{email}{mailing address}
%\cvref{Prof.\ Alpha Beta}{Institute}{a.beta@university.edu}
%{Address Line 1\\Address line 2}

%\divider



\end{paracol}

\end{document}
